\chapter{The Effective Potential in Cosmology}
\label{ch:eff-pot-cosmo}

In Part~\ref{part:eff-action} we developed the effective action
formalism and computed the effective potential at zero temperature,
using a momentum cutoff to regulate the divergent integrals. Here we
revisit this computation in the cosmological context, with two
differences. First, we use dimensional regularization and the
$\overline{\text{MS}}$ renormalization scheme, which is more systematic
and avoids the ambiguities associated with a hard cutoff. Second, we
keep the physical application front and center: the question is not just
what the quantum-corrected vacuum is, but how it changes as the universe
cools, and what determines whether the resulting phase transition is
first or second order.

\section{The Effective Action and Effective Potential}

We work with a real scalar field $\phi$ in $\phi^4$ theory,
\begin{equation}
	S[\phi] = \int d^4x\left[
		\frac{1}{2}\partial_\mu\phi\,\partial^\mu\phi
		- V_0(\phi)\right],\qquad
	V_0(\phi) = \frac{m^2}{2}\phi^2 + \frac{\lambda}{4!}\phi^4,
\end{equation}
where the $\phi\to-\phi$ symmetry forbids a cubic term. The basic
objects of interest are the correlators
\begin{equation}
	\langle\Omega|T\{\phi(x_1)\cdots\phi(x_n)\}|\Omega\rangle
	= \frac{\int\mathcal{D}\phi\;\phi(x_1)\cdots\phi(x_n)\,
		e^{iS[\phi]}}{\int\mathcal{D}\phi\;e^{iS[\phi]}},
\end{equation}
where $|\Omega\rangle$ is the vacuum of the interacting theory. Our goal
is to find a potential for the expectation value
$\langle\Omega|\phi(x)|\Omega\rangle$.

Introducing a source $J(x)$, the generating functional is
\begin{equation}
	Z[J] = \int\mathcal{D}\phi\;
	e^{iS[\phi] + i\int d^4x\,J(x)\phi(x)},
\end{equation}
from which all correlators are obtained by functional differentiation.
The connected generating functional $W[J] = -i\ln Z[J]$ generates the
connected Green's functions. Its first derivative defines the classical
field,
\begin{equation}
	\frac{\delta W[J]}{\delta J(x)}
	= \langle\Omega|\phi(x)|\Omega\rangle_J
	\equiv \bar\phi(x).
\end{equation}
The effective action is defined by the Legendre transform,
\begin{equation}
	\Gamma[\bar\phi] \equiv W[J]
	- \int d^4x\,J(x)\,\bar\phi(x),
\end{equation}
and satisfies $\delta\Gamma/\delta\bar\phi(y) = -J(y)$. The vacuum of
the theory in the absence of sources is therefore determined by
\begin{equation}
	\frac{\delta\Gamma[\bar\phi]}{\delta\bar\phi} = 0.
\end{equation}
For a spatially uniform, time-independent configuration
$\bar\phi(x) = \bar\phi_c$, the effective action reduces to
\begin{equation}
	\Gamma[\bar\phi_c] = -\int d^4x\,V_{\text{eff}}(\bar\phi_c),
\end{equation}
where $V_{\text{eff}}(\bar\phi_c)$ is the effective potential, the
quantum-corrected potential whose minimum determines the true ground
state of the theory. Without proof, $\Gamma[\bar\phi]$ is the
generating functional of 1PI correlators,
\begin{equation}
	V_{\text{eff}}(\bar\phi_c) = -\sum_{n=0}^{\infty}
	\frac{1}{n!}\,\bar\phi_c^{\,n}\,\Gamma^{(n)}(p_i=0).
\end{equation}

\section{The Coleman-Weinberg Potential via Dimensional Regularization}

The one-loop correction to the effective potential is
\begin{equation}
	V_1(\phi) = \frac{1}{2}\int\frac{d^4p_E}{(2\pi)^4}
	\ln\!\left[p_E^2 + \tilde m^2(\phi)\right],
\end{equation}
where
\begin{equation}
	\tilde m^2(\phi) = \frac{d^2 V_0(\phi)}{d\phi^2}
	= m^2 + \frac{\lambda}{2}\phi^2
\end{equation}
is the field-dependent mass. This integral is UV divergent and must be
regulated. We use dimensional regularization, replacing $d=4$ by
$d=4-2\varepsilon$ and introducing a mass scale $\mu$ to keep the
action dimensionless,
\begin{equation}
	V_1(\phi) = \frac{1}{2}\mu^{4-d}
	\int\frac{d^dk_E}{(2\pi)^d}
	\ln\!\left(p_E^2 + \tilde m^2(\phi)\right).
\end{equation}
To evaluate this, we differentiate with respect to $\tilde m^2(\phi)$
and use the master formula of dimensional regularization,
\begin{equation}
	\int\frac{d^dp_E}{(2\pi)^d}
	\frac{1}{p_E^2+m^2}
	= \frac{1}{(4\pi)^{d/2}}(m^2)^{d/2-1}
	\Gamma\!\left(1-\frac{d}{2}\right),
\end{equation}
to obtain
\begin{equation}
	\frac{dV_1}{d\tilde m^2}
	= \frac{1}{2}\mu^{4-d}\cdot
	\frac{1}{(4\pi)^{d/2}}
	\left(\tilde m^2\right)^{d/2-1}
	\Gamma\!\left(1-\frac{d}{2}\right).
\end{equation}
Integrating back and substituting $d=4-2\varepsilon$,
\begin{equation}
	V_1(\phi) = \frac{1}{64\pi^2}\,\tilde m^4(\phi)
	\left[-\frac{1}{\varepsilon} + \gamma_E - \ln 4\pi
		+ \ln\frac{\tilde m^2(\phi)}{\mu^2} - \frac{3}{2}\right].
\end{equation}
In the $\overline{\text{MS}}$ scheme, we absorb the divergent combination
$\mathcal{C}_{\text{uv}} = 1/\varepsilon - \gamma_E + \ln 4\pi$ into
counterterms, leaving the renormalized one-loop potential
\begin{equation}
	\boxed{V_1(\phi) = \frac{1}{64\pi^2}\,\tilde m^4(\phi)
		\left[\ln\frac{\tilde m^2(\phi)}{\mu^2} - \frac{3}{2}\right].}
	\label{eq:CW-msbar}
\end{equation}
This result generalizes immediately to other fields. For a $\mathrm{U}(1)$
gauge field with field-dependent mass $\tilde m_A^2(\phi) = 2g^2|\phi|^2$
arising from spontaneous symmetry breaking,
\begin{equation}
	V_1^{(A)}(\phi) = \frac{3}{64\pi^2}\,\tilde m_A^4(\phi)
	\left[\ln\frac{\tilde m_A^2(\phi)}{\mu^2} - \frac{5}{6}\right],
\end{equation}
where the factor of 3 counts the physical polarizations and the
$-5/6$ (rather than $-3/2$) arises from the vector nature of the field,
computed in Landau gauge. For Dirac fermions with Yukawa coupling $y$
and field-dependent mass $m_\psi^2 = y^2\phi^2$,
\begin{equation}
	V_1^{(\psi)}(\phi) = -\frac{4}{64\pi^2}\,\tilde m_\psi^4(\phi)
	\left[\ln\frac{\tilde m_\psi^2(\phi)}{\mu^2} - \frac{3}{2}\right],
\end{equation}
where the overall minus sign comes from the fermion loop and the factor
of 4 counts the spinor degrees of freedom. The general formula,
summing over all particles coupling to $\phi$, is
\begin{equation}
	V_{\text{eff}}(\phi) = V_0(\phi)
	+ \frac{1}{64\pi^2}\sum_i n_i\,\tilde m_i^4(\phi)
	\left[\ln\frac{\tilde m_i^2(\phi)}{\mu^2} - c_i\right],
	\label{eq:CW-general}
\end{equation}
where $n_i = (-1)^F\cdot\text{d.o.f.}$ counts degrees of freedom with
a minus sign for fermions, and
\begin{equation}
	c_i = \begin{cases}
		\dfrac{3}{2} & \text{scalars and fermions,} \\[6pt]
		\dfrac{5}{6} & \text{gauge bosons.}
	\end{cases}
\end{equation}
This is the Coleman-Weinberg potential in the $\overline{\text{MS}}$
scheme. It supersedes the cutoff-regulated result of
Chapter~\ref{ch:toy-model} and will be our working formula for
the remainder of this part.

\section{The Coleman-Weinberg Model}

We now apply \eqref{eq:CW-general} to the Coleman-Weinberg model of
massless scalar electrodynamics,
\begin{equation}
	\mathcal{L} = -\frac{1}{4}F_{\mu\nu}^2
	+ |D_\mu\phi|^2 - \frac{\lambda}{4!}|\phi|^4,
	\qquad D_\mu = \partial_\mu + igA_\mu.
\end{equation}
There is no mass term in the Lagrangian, and the classical action is
scale invariant: under $x\to\lambda x$, $A_\mu\to\lambda^{-1}A_\mu$,
$\phi\to\lambda^{-1}\phi$, each term in $\mathcal{L}$ scales as
$\lambda^{-4}$, leaving $S = \int d^4x\,\mathcal{L}$ invariant.

At the quantum level, scale invariance is broken by the logarithms in
\eqref{eq:CW-general}. Writing $\phi = \phi_c + \chi_1 + i\chi_2$ and
expanding around the background $\phi_c$, the field-dependent masses are
\begin{align}
	\tilde m_{\chi_1}^2(\phi_c) & = \frac{\lambda}{2}\phi_c^2, \\
	\tilde m_{\chi_2}^2(\phi_c) & = \frac{\lambda}{6}\phi_c^2, \\
	m_A^2(\phi_c)               & = 2g^2\phi_c^2.
\end{align}
In the regime $\lambda^2\ll g^4$, the scalar loop contributions are
negligible and the effective potential is dominated by the gauge boson
loop,
\begin{equation}
	V_{\text{eff}}(\phi_c) \cong \frac{\lambda}{4!}\phi_c^4
	+ \frac{3}{16\pi^2}\,g^4\phi_c^4
	\left[\ln\frac{2g^2\phi_c^2}{\mu^2} - \frac{5}{6}\right].
\end{equation}
Redefining the renormalization scale as $\tilde\mu^2 =
	e^{5/6}\mu^2/(2g^2)$, this simplifies to
\begin{equation}
	V_{\text{eff}}(\phi_c) = \frac{\lambda}{4!}\phi_c^4
	+ \frac{3g^4}{16\pi^2}\,\phi_c^4
	\ln\frac{\phi_c^2}{\tilde\mu^2}.
\end{equation}
Defining an effective quartic coupling
\begin{equation}
	\lambda_{\text{eff}}(\phi_c) \equiv \lambda
	+ \frac{g^4}{2\pi^2}\ln\frac{\phi_c^2}{\tilde\mu^2},
\end{equation}
the effective potential takes the suggestive form
\begin{equation}
	V_{\text{eff}}(\phi_c) = \frac{\lambda_{\text{eff}}(\phi_c)}{4!}
	\phi_c^4.
\end{equation}
This is precisely analogous to the running coupling in scattering
experiments, where the appropriate scale at which to evaluate the
coupling is the energy scale of the process. Here the relevant scale is
the field value $\phi_c$ itself. At $\phi_c = \tilde\mu$, the logarithm
vanishes and $\lambda_{\text{eff}}(\tilde\mu) = \lambda$, so $\tilde\mu$
plays the role of an RG scale. Quantum corrections have therefore broken
the classical scale invariance and generated a preferred scale, producing
a nontrivial minimum of the effective potential.

% [FIGURE: Coleman-Weinberg effective potential showing dynamically
%  generated minimum]

This is the phenomenon of dimensional transmutation already encountered
in Chapter~\ref{ch:massless-scalar-ed}: a theory with no classical mass
scale develops a preferred scale through radiative corrections. The scale
$\tilde\mu$ at which $\lambda_{\text{eff}}$ is evaluated plays the same
role as the renormalization point $M$ in the cutoff treatment of
Chapter~\ref{ch:toy-model}, and the two approaches give the same
physical result in different renormalization schemes.

For the Standard Model Higgs, an analogous but more involved running
occurs. The tree-level wine-bottle potential has a minimum at
$\sim 100\,\text{GeV}$ with positive $\lambda$. However, running $\lambda$
to higher scales using the Standard Model renormalization group equations,
the coupling turns negative at energy scales of order $10^{10}$ to
$10^{12}\,\text{GeV}$, implying that the electroweak vacuum is
metastable rather than absolutely stable.

% [FIGURE: Higgs effective potential showing metastable minimum]

The lifetime of this metastable vacuum is many orders of magnitude
larger than the age of the universe, so this poses no immediate
phenomenological problem. Nevertheless it raises the deep question of
why the measured values of the top quark mass and Higgs mass place us
precisely in the metastable region of parameter space. For a detailed
analysis, see Buttazzo et al., arXiv:1307.3536.
